\section{Comparison with covariance-based allocators}
\label{sec:comparison}

Minimum variance, hierarchical risk parity and their relatives are functions of a
covariance matrix. The race differs from them in two respects. It never inverts
the covariance, and it can respond to tail dependence.

\subsection{Avoiding matrix inversion}
\label{sec:robust}

The operational fragility of mean--variance optimization is the inversion of an
estimated covariance \parencite{markowitz1952}: as $\Sigma$ nears singularity,
$\Sigma^{-1}$ and with it the weights explode. The discipline's responses are
shrinkage toward a structured target \parencite{ledoitwolf2004} and, more
radically, allocators that never form $\Sigma^{-1}$ at all (hierarchical risk
parity \parencite{lopezdeprado2016}, the Schur-complementary family
\parencite{cotton2024schur}, risk parity \parencite{maillard2010}), trading point
optimality for stability and scale. The Thurstone race belongs to this inversion-free family: it
\emph{samples} from a correlation rather than inverting one, so it stays well posed
even when $\Sigma$ is singular or rank-deficient.

Covariance-based allocators also split a position between near-duplicates, but
unstably. Mean--variance settles the split by inverting $\Sigma$, and the direction
in which two assets become identical is exactly the direction in which $\Sigma$
loses rank; so the population answer is right while its sample estimate, dividing
through a vanishing eigenvalue, swings between the twins and flips sign with the
data \parencite{markowitz1952}.

The clustering allocators, such as hierarchical risk parity
\parencite{lopezdeprado2016}, and the Schur-complementary family
\parencite{cotton2024schur} that relates these to the minimum variance portfolio,
de-duplicate by grouping instead, but discretely: the split jumps as cluster
membership flips. The race de-duplicates \emph{smoothly}, with a closed
characterization (\Cref{lem:redundancy}) and without ever forming $\Sigma^{-1}$.

Redundancy is near-singularity, so robustness to ill-conditioning, low turnover
(\Cref{thm:smooth}, \Cref{sec:experiments}) and the graceful split of redundant
exposure are one behaviour. The weights move smoothly along the direction in which
assets become collinear and a covariance inverse would explode.

Robustness to \emph{conditioning} is not robustness to \emph{tails}. Every method
just named is a functional of the covariance, or of a correlation derived from it.
A covariance is a second moment; it summarises how assets move together on average
and says nothing about how they move together in the extreme. Two universes with
identical covariance, hence identical minimum-variance, HRP, Schur, and
risk-parity weights, can carry entirely different probabilities of a simultaneous
crash. The quantity that separates them, tail dependence, is not a moment, and no
re-estimation or shrinkage of $\Sigma$ recovers it.

\subsection{Empirical motivation}
\label{sec:tailmotiv}

That equity returns are more dependent in falling markets than a Gaussian model
allows is well documented. \textcite{longinsolnik2001} reject normality for the
joint lower tail of international equity returns but not for the upper tail, and
\textcite{angchen2002} find that correlations between U.S. equity portfolios and
the market are higher in downturns than a normal distribution implies.
\textcite{embrechts2002} set out why correlation is an incomplete summary of
dependence. \textcite{poon2004} distinguish asymptotic dependence from asymptotic
independence in the joint tails of stock indices and show that traditional
dependence measures can misstate portfolio risk, and \textcite{patton2004} finds
economically significant gains from modelling asymmetric dependence in an
asset-allocation problem. A standard summary is the lower-tail dependence
coefficient
\begin{equation}
  \lambda_L \;=\; \lim_{u\downarrow 0}\Prob\!\big(U_i < u \,\big|\, U_j < u\big)
  \label{eq:lambdaL}
\end{equation}
with $U_i = F_i(R_i)$ the probability-integral transform of asset $i$'s return,
which is zero for every Gaussian copula with correlation below one
\parencite{sibuya1960}.

We add to this existing motivation a simple observation about the Dow 30. A
multivariate Gaussian fit to the daily returns of its constituents, pinning the
means and the full covariance, reproduces a shallow joint decline well enough, but
its joint tail thins far too fast. All of them closing down on the same day is only
mildly underpredicted; deepen the event and the gap widens without bound.

Over 2010--2024, days on which every name fell by more than two percent occurred
about eighty times more often than the fitted Gaussian allows, and beyond three
percent the Gaussian effectively never produces such a day, yet two did. The same
fit underpredicts joint \emph{crashes} far more than joint \emph{rallies}
(\Cref{tab:tail}).\footnote{Daily log returns of
the 28 Dow names with full 2010--2024 history (Yahoo Finance; some $3{,}770$ trading days; average pairwise
correlation $0.42$), against a Gaussian fit to their mean and covariance, the
probabilities by $10^{8}$ Monte-Carlo draws. The deepest rows rest on two to four
days, so the precise multiples are noisy; the monotone, asymmetric divergence is
not.}

\begin{table}[ht]
  \centering
  \caption{All $28$ Dow constituents moving by more than $q$ on the same day
  (2010--2024): empirical frequency, the frequency under a Gaussian fit to the
  mean and covariance, and their ratio. The Gaussian undercounts joint
  \emph{crashes} ever more steeply with depth, and far more than joint
  \emph{rallies}: the lower-tail dependence a second moment cannot carry.}
  \label{tab:tail}
  \begin{tabular}{lccc@{\qquad}ccc}
    \toprule
    & \multicolumn{3}{c}{all down by ${>}q$} & \multicolumn{3}{c}{all up by ${>}q$} \\
    \cmidrule(lr){2-4}\cmidrule(lr){5-7}
    $q$ & emp. & Gaussian & ratio & emp. & Gaussian & ratio \\
    \midrule
    $1.0\%$ & $0.32\%$ & $0.080\%$  & $4\times$         & $0.24\%$ & $0.111\%$  & $2\times$  \\
    $1.5\%$ & $0.21\%$ & $0.012\%$  & $18\times$        & $0.13\%$ & $0.018\%$  & $7\times$  \\
    $2.0\%$ & $0.11\%$ & $0.0014\%$ & $77\times$        & $0.05\%$ & $0.0022\%$ & $25\times$ \\
    $3.0\%$ & $0.05\%$ & ${\sim}0$  & ${>}10^{4}\times$ & $0\%$    & ${\sim}0$  & n/a        \\
    \bottomrule
  \end{tabular}
\end{table}

The table changes the marginals and the copula together, so by itself it shows
only that the moment-matched Gaussian fails. Keeping each name's empirical marginal
and making only the copula Gaussian (normal-scores correlation, $10^8$ draws) does
worse still: all down by more than $2\%$ has probability $5\times10^{-6}$, against
$1.4\times10^{-5}$ for the full Gaussian fit and $0.11\%$ observed. The shortfall
lies in the dependence, not the marginals.

A race among assets reads the joint extreme directly. The winner of a race is an
\emph{order statistic} of the joint draw, and the law of an order statistic is set
by the tail copula, not by the covariance. Drive the race with a
simulation that carries lower-tail dependence and the weights respond. What governs
the response is the distribution of a cluster's worst performance relative to the
rest of the field, not the pairwise coefficient $\lambda_L$ alone
(\Cref{lem:tail}). No covariance-based allocator can respond at all.

The race sees dependence only through relative performances, so a shock common to
every asset leaves the weights unchanged. A response requires that losses fall on
some names relative to others, as when a cluster crashes together and the rest of
the field does not.


\section{The race as a choice model}
\label{sec:choicemodel}

When the present methodology is applied to tilt capitalization-weighted indexes,
the contrast has an interesting parallel to old debates in psychometry and
preference. The race is Thurstone's model of choice. Capitalization weighting
corresponds to the rival model, Luce's, and the two differ in how weights respond
when the field of assets changes.

\subsection{Asset allocation as a choice}
\label{sec:choice}

The race of \Cref{sec:method} is the law of comparative judgment of
\textcite{thurstone1927}, the canonical model of probabilistic choice. Each
alternative carries a latent ability and a noisy performance, and the one with the
extreme performance is chosen. On this reading the market selects among assets, and
a portfolio weight is the probability that an asset is chosen.

The tilt is informationally modest. It needs only the dependence, no return
forecast, and it \emph{anchors} to the benchmark, departing from the prevailing
allocation only as far as the correlation warrants.

\subsection{Luce's axiom and capitalization weighting}
\label{sec:luce}

The choice literature has a rival account, Luce's \parencite{luce1959}, whose
multinomial form is the logit. Write $w_i(\mathcal U)$ for the weight a rule assigns to asset $i$ within a
universe $\mathcal U$. Capitalization weighting sets $w_i(\mathcal U) =
\mathrm{cap}_i / \sum_{j\in\mathcal U}\mathrm{cap}_j$, so the ratio
$w_i(\mathcal U)/w_j(\mathcal U) = \mathrm{cap}_i/\mathrm{cap}_j$ does not depend on
$\mathcal U$. That is Luce's choice axiom, independence of irrelevant
alternatives, and we call such a rule \emph{Lucian}.

Imagine you are a psychometrician trying to understand the market's choices, and
how they change when assets are removed. Luce's axiom is one hypothesis about that.
It would be reasonable to test a range of choice models against the market's
behaviour, not only Luce's. In a Thurstonian race the response to a removal depends
on how the removed asset co-moves with the rest (\Cref{lem:redundancy}).

The Thurstonian family is the older and the more general of the
two.\footnote{Luce's model is equivalent to a race with Gumbel performances
\parencite{yellott1977}.} It lost the applied field to the logit on tractability:
beyond two alternatives its Gaussian member, the probit, has no closed form, and
recovering abilities from choice shares was considered prohibitive. The inversion
of \Cref{sec:method} removes that obstacle.

\subsection{Why the CAPM does not rescue naive indexing}
\label{sec:capm}

The passive industry defends capitalization weighting by appeal to the CAPM. We
grant the CAPM in full. The standard argument \parencite{sharpe1964} is that every
investor is a mean--variance optimizer who agrees on the expected returns $\mu$ and
on the same covariance $\Sigma$, so each holds the risk-free asset together with the
same tangency portfolio, proportional to $\Sigma^{-1}(\mu - r\mathbf{1})$. Market
clearing forces that portfolio to equal the supply, the capitalization-weighted
market, so $\mathrm{cap} \propto \Sigma^{-1}(\mu - r\mathbf{1})$ and the market
portfolio is efficient.\footnote{One wrinkle we do not dwell on:
investors have a well-established motivation to shrink the covariance
\parencite{ledoitwolf2004}, and an aggregate of heterogeneous shrunk estimates
need not clear to a mean--variance-efficient market portfolio, so even the
equilibrium claim is not airtight.}

That endorsement is for the whole-market universe and does not transfer to the
restricted universes where weighting decisions are made. This is Roll's critique
\parencite{roll1977}: the only testable content of the CAPM is the efficiency of
the \emph{true} market portfolio, and no proxy on a restricted universe inherits it
\parencite{prono2009}. The block calculation below states the point in the present
notation. We do not claim it as new.

\begin{lemma}[Lucian weighting is inconsistent under restriction]
\label{lem:restriction}
Capitalization weighting is the Lucian rule with values $\mathrm{cap}_i$. If
returns are correlated (non-diagonal $\Sigma$), then generically a Lucian rule
coincides with mean--variance efficiency on at most one universe in a nested
family. Under the CAPM that universe is the full asset world $U$; on any strict
sub-universe $S \subsetneq U$ the capitalization-weighted portfolio is generically
inefficient, and the discrepancy is governed by the heterogeneity of the assets'
covariances with the excluded set $U \setminus S$.
\end{lemma}
\begin{proof}
Partition the market $U = S \cup X$ into the retained assets $S$ and the excluded
$X = U\setminus S$, with market weights $m = (m_S, m_X)$. CAPM equilibrium reads
$\mu - r\mathbf 1 = \gamma\,\Sigma_U\, m$ for some $\gamma > 0$; its $S$-block is
$\mu_S - r\mathbf 1 = \gamma(\Sigma_{SS} m_S + \Sigma_{SX} m_X)$, so the
restricted-universe tangency weights are
\begin{equation*}
  q_S \;\propto\; \Sigma_{SS}^{-1}(\mu_S - r\mathbf 1)
  \;=\; \gamma\big(m_S + \Sigma_{SS}^{-1}\Sigma_{SX}\, m_X\big),
\end{equation*}
while restricted capitalization weighting uses $m_S$. The two coincide iff
$\Sigma_{SS}^{-1}\Sigma_{SX}\, m_X = c\, m_S$ for some scalar $c > -1$, the sign
keeping the normalised weights on the simplex. At the full universe
$X = \varnothing$ and $c = 0$. On a strict sub-universe the condition is $|S|$
equations in the one unknown $c$, which fail for generic $\Sigma_{SX} \neq 0$.

The non-generic exceptions are a block-diagonal $\Sigma$ with $\Sigma_{SX} = 0$,
and an excluded set whose market portfolio regresses identically on every retained
asset, $\Sigma_{SX} m_X \propto \Sigma_{SS} m_S$. The wedge
$\Sigma_{SS}^{-1}\Sigma_{SX}\, m_X$, the coefficients from regressing the return of
the excluded portfolio $m_X$ on the retained assets' returns, is the part of each cap weight that reflects covariance
with $U\setminus S$, irrelevant to the $S$-frontier yet baked into the full-market
weight.
\end{proof}

Fundamental indexing reaches cap weighting's sub-optimality by a different route,
price inefficiency \parencite{hsu2006}, while shrinkage \parencite{ledoitwolf2004}
and tracking-error analysis \parencite{roll1992} document the fragility of the
inputs.
